% Chapter 3 -- Claim: Individual Investors Hurt Themselves
% Requires: amsmath, amssymb, booktabs, array, tabularx, graphicx, tcolorbox (breakable, skins), hyperref or url.
% Defines the callout box only if the preamble has not.
\providecolor{fenavy}{HTML}{1B2A6B}
\providecolor{feshade}{HTML}{F3F5FA}
\makeatletter
\@ifundefined{callout}{%
  \newtcolorbox{callout}{breakable, enhanced,
    colback=white, colframe=gray!70, boxrule=0.5pt, arc=1mm,
    borderline west={2.5pt}{0pt}{gray!70},
    left=8pt, right=8pt, top=4pt, bottom=4pt}%
}{}
\makeatother

\chapter{Claim: Individual Investors Hurt Themselves}
\label{ch:03}

Classical finance treats the investor as a rational agent who holds a diversified portfolio, trades only when new information arrives, and earns the return the market offers. Starting in the late 1990s, researchers gained access to something the theory had never been tested against: the actual trading records of ordinary households. Odean (1998) obtained the accounts of 10,000 customers of a discount broker, and Barber and Odean (2000) extended the sample to 66,465 households. For the first time it was possible to ask not whether markets are efficient, but whether the people who trade in them behave as the theory assumes.

The answer mattered for two reasons. For research, systematic departures from rational trading would support the new field of behavioral finance, which argued that psychology, not only risk, shapes prices and returns. For practice, it would mean that the return an investor earns can differ from the return an investment earns, and that the difference is under the investor's own control. Dichev (2007) carried the question from individual accounts to whole markets. He showed that the dollar-weighted return of the aggregate investor, which counts when money arrives and leaves, falls short of the buy-and-hold return, because capital tends to flow in before poor returns and out before good ones. Frazzini and Lamont (2008) found the same pattern in mutual fund flows and called it ``dumb money''.

The implications reach far beyond the academic debate. If investors lose return through their own decisions, then product design, default options, financial advice and regulation can raise their wealth without finding a single new source of return. The same evidence also shows why some popular figures overstate the problem, because part of any gap reflects when people have money to save, not mistakes. The claim on trial\index{claims-on-trial method} is therefore: \textbf{individual investors lose a meaningful part of the return on their investments through their own trading and timing decisions.}

The chapter draws on four papers. Odean (1998) measures the disposition effect, Barber and Odean (2000) measure the cost of overtrading, Dichev (2007) measures the gap between dollar-weighted and buy-and-hold returns, and Frazzini and Lamont (2008) link fund flows to later stock returns. Because individual brokerage records are private, the lab in Section~\ref{sec:03-lab} replicates Dichev's market-level calculation exactly with public Federal Reserve data, measures the same gap for mutual fund investors with public fund-flow data, and reproduces the logic of Odean and of Barber and Odean with simulated accounts trading real industry portfolios. It does this in order to measure how much return each behavior costs, and how much of the measured gap reflects the timing of savings rather than mistakes.

By the end of the chapter you will be able to:

\begin{itemize}
  \item compute time-weighted and money-weighted returns\index{return!money-weighted}, and the behavior gap\index{behavior gap} between them;
  \item compute the dollar-weighted return of a whole stock market from its net equity issuance and dividends, as Dichev (2007) did;
  \item measure whether aggregate fund flows chase past returns, and whether they predict future returns;
  \item name the main behavioral biases of individual investors: overconfidence\index{overconfidence}, the disposition effect\index{disposition effect}, performance chasing\index{performance chasing}, attention, home bias\index{home bias} and lottery preference\index{lottery preference};
  \item measure the disposition effect with the proportions of gains and losses realized;
  \item explain sequence-of-returns risk\index{sequence-of-returns risk} and the 4\% withdrawal rule;
  \item state the three parts of the claim and judge the evidence, including why some widely quoted figures overstate the gap;
  \item design rules (automatic contributions, calendar rebalancing\index{rebalancing}, a written policy) that take decisions out of the moment.
\end{itemize}

\section{Foundations: measuring investor returns, behavioral biases and withdrawals}
\label{sec:03-foundations}

\subsection{Time-weighted return}

The \textbf{time-weighted return}\index{return!time-weighted|textbf} (TWR) measures what one dollar invested at the start and left alone would have earned. It links the period returns and ignores money added or withdrawn:
\begin{equation}
  1 + \mathrm{TWR} = (1 + r_1)(1 + r_2)\cdots(1 + r_T)
  \label{eq:03-twr}
\end{equation}
Fund performance is reported this way, because the manager does not control when investors add or remove money.

\subsection{Money-weighted return}

The \textbf{money-weighted return}\index{return!money-weighted|textbf} (MWR), or dollar-weighted return, is the internal rate of return of the investor's actual cash flows: the rate $y$ that sets
\begin{equation}
  -V_0 - \sum_{t=1}^{T-1}\frac{C_t}{(1 + y)^t} + \frac{V_T}{(1 + y)^T} = 0
  \label{eq:03-mwr}
\end{equation}
where $V_0$ is the starting value, $C_t$ the contributions (negative for withdrawals) and $V_T$ the ending value. It weights each period by the money invested at the time, so it measures what the investor actually earned.

\subsection{The behavior gap}

The difference MWR $-$ TWR is often called the \textbf{behavior gap}\index{behavior gap|textbf} or investor return gap. It is negative when money tends to arrive before poor returns and leave before good ones.

\textbf{Worked example.} A fund gains 20\% in year 1 and loses 10\% in year 2 (Figure~\ref{fig:03-twr}).

\begin{itemize}
  \item Its TWR is $1.20 \times 0.90 - 1 = 8\%$ over two years, about 3.9\% a year.
  \item An investor puts in \$1,000 at the start and, impressed by the first year, adds \$1,000 at the start of year 2. After year 1 the account holds \$1,200, then \$2,200 after the deposit, and \$1,980 at the end.
  \item Solving $-1{,}000 - 1{,}000/(1 + y) + 1{,}980/(1 + y)^2 = 0$ gives $y \approx -0.7\%$ a year.
  \item The fund earned 3.9\% a year, but the investor lost money, because more dollars were exposed to the bad year than to the good one.
\end{itemize}

\begin{figure}[htbp]
\centering
\scalebox{0.8}{%
\begin{tikzpicture}[x=1cm, y=0.0021cm]
  % Bars to scale: dollars at work in each year
  \draw (-0.3,0) -- (5.3,0);
  \fill[fenavy!35] (0.2,0) rectangle (2.0,1000);
  \fill[fenavy]    (2.9,0) rectangle (4.7,2200);
  \node[anchor=south] at (1.1,1000) {\$1,000 at work};
  \node[anchor=south] at (3.8,2200) {\$2,200 at work};
  \node[anchor=north] at (1.1,-40) {Year 1: $+20\%$};
  \node[anchor=north] at (3.8,-40) {Year 2: $-10\%$};
  % Results
  \node[anchor=west, align=left] at (5.8,1500) {fund (time-weighted): $+3.9\%$ a year};
  \node[anchor=west, align=left] at (5.8,1120) {investor (money-weighted): $-0.7\%$ a year};
  \node[anchor=west, align=left, fegray] at (5.8,680)
    {more dollars were exposed\\to the bad year than the good one};
\end{tikzpicture}%
}
\caption{The worked example. The fund's good year was earned on \$1,000 and its bad year on \$2,200, so the investor's money-weighted return is $-0.7\%$ a year while the fund reports $+3.9\%$.}
\label{fig:03-twr}
\end{figure}

\subsection{Aggregate flows and the dollar-weighted market return}
\label{sec:03-flows}

The same calculation can be applied to a whole market. Treat all shareholders together as one investor. Each year this investor receives dividends and the cash paid out in share repurchases, and pays in the cash raised by new share issues. The net cash flow to the investor is
\begin{equation}
  D_t = \text{dividends}_t - \text{net equity issues}_t
  \label{eq:03-netflow}
\end{equation}
where net equity issues are new issues minus repurchases, including the shares retired in cash takeovers. When companies issue more than they pay out, $D_t$ is negative: investors as a group are putting money in.

Dichev (2007) computed the \textbf{dollar-weighted market return}\index{return!dollar-weighted market|textbf} as the internal rate of return that equates the market value at the start, the net flows $D_t$ in between, and the market value at the end. It is Equation~\eqref{eq:03-mwr} applied to the market. Comparing it with the buy-and-hold return, which ignores the flows, shows whether money tends to arrive at good or bad times.

\textbf{Fund flows}\index{fund flows|textbf} are the same idea one level down. The \textbf{net new cash flow} into a group of mutual funds is sales minus redemptions, plus net exchanges from other funds. Flows are usually scaled by the funds' assets at the start of the month, so that a flow of 1\% means investors added a hundredth of the money already invested. Performance chasing shows up as a positive relation between this month's flow and past returns; ``dumb money'' shows up as a negative relation between flows and future returns.

\subsection{Behavioral biases}

Behavioral finance\index{behavioral finance}, which Chapter~\ref{ch:05} uses to explain momentum\index{momentum}, documents systematic departures by individual investors from the textbook investor (Barberis and Thaler 2003):

\begin{itemize}
  \item \textbf{Overconfidence}\index{overconfidence|textbf} leads investors to trade too much, believing they can pick winners.
  \item \textbf{The disposition effect}\index{disposition effect|textbf} is the tendency to sell winners too soon and hold losers too long (Shefrin and Statman 1985). \textbf{Prospect theory}\index{prospect theory|textbf} offers one explanation: investors feel losses more than equal gains and dislike realizing them.
  \item \textbf{Extrapolation and performance chasing}: buying what has recently risen, in the belief that it will keep rising.
  \item \textbf{Attention}: buying stocks that are in the news or top the trading-app lists.
  \item \textbf{Home bias and under-diversification\index{under-diversification}}: holding too few stocks, often from one's own country, region or employer.
  \item \textbf{Lottery preference}\index{lottery preference|textbf}: paying up for stocks with a small chance of a very large gain (Chapter~\ref{ch:18}).
\end{itemize}

\subsection{Measuring the disposition effect}

Odean (1998) compared, on each day an investor sold something, the \textbf{proportion of gains realized}\index{disposition effect!PGR and PLR|textbf} (PGR) with the \textbf{proportion of losses realized} (PLR):
\begin{equation}
  \mathrm{PGR} = \frac{\text{realized gains}}{\text{realized gains} + \text{paper gains}}, \qquad
  \mathrm{PLR} = \frac{\text{realized losses}}{\text{realized losses} + \text{paper losses}}
  \label{eq:03-pgr}
\end{equation}
A PGR above the PLR means investors sell winners more readily than losers.

\subsection{Commissions and spreads in brief}
\label{sec:03-costs}

Every trade has a cost, even when the broker charges no commission. A buyer pays the \textbf{ask}, the price at which others are willing to sell, and a seller receives the \textbf{bid}, the price at which others are willing to buy. The gap between them is the \textbf{bid--ask spread}\index{bid--ask spread}, and half of it is the cost of trading in one direction. For a portfolio, the yearly cost is roughly the cost of a round trip times the share of the portfolio replaced each year, its \textbf{turnover}\index{turnover}. At a round-trip cost of 1\% and a turnover of 75\%, a portfolio loses about 0.75 percentage points a year. Chapter~\ref{ch:17} measures these costs in detail; this chapter needs only the arithmetic.

\subsection{Withdrawal rates}

A retiree living off a portfolio faces \textbf{sequence-of-returns risk}\index{sequence-of-returns risk|textbf}: the same average return is far more damaging if the bad years come early, when withdrawals take a larger share of a shrinking portfolio (Figure~\ref{fig:03-sequence}). Bengen (1994) asked how large a fixed, inflation\index{inflation}-adjusted withdrawal would have survived every 30-year period in U.S. history, and found about 4\% of the initial portfolio.

\begin{figure}[htbp]
\centering
\begin{tikzpicture}
\begin{axis}[
  width=0.64\linewidth, height=0.34\linewidth,
  axis lines=left, axis line style={-stealth},
  xmin=0, xmax=30, ymin=0, ymax=3200,
  xtick={0,5,...,30}, ytick={0,1000,2000,3000},
  yticklabel style={/pgf/number format/1000 sep={,}},
  scaled y ticks=false,
  xlabel={Years into retirement},
  ylabel={Portfolio value (\$ thousands, real)},
  clip=false,
]
  % $1,000k start, $40k withdrawn at the start of each year;
  % returns: three years of -15% and twenty-seven of +7%, in two orders
  \addplot[fegray, very thick, dashed] coordinates {
    (0,1000) (1,816) (2,660) (3,527) (4,521) (5,514) (6,508) (7,500) (8,493)
    (9,484) (10,475) (11,466) (12,456) (13,445) (14,433) (15,421) (16,407)
    (17,393) (18,378) (19,361) (20,344) (21,325) (22,305) (23,283) (24,260)
    (25,236) (26,210) (27,181) (28,151) (29,119) (30,85)};
  \addplot[fenavy, very thick] coordinates {
    (0,1000) (1,1027) (2,1056) (3,1087) (4,1121) (5,1156) (6,1195) (7,1235)
    (8,1279) (9,1326) (10,1376) (11,1429) (12,1487) (13,1548) (14,1613)
    (15,1684) (16,1759) (17,1839) (18,1925) (19,2017) (20,2115) (21,2220)
    (22,2333) (23,2453) (24,2582) (25,2720) (26,2868) (27,3026) (28,2538)
    (29,2123) (30,1771)};
  \node[fenavy, align=center] at (axis cs:10,2350) {bad years last:\\ends at \$1.77 million};
  \node[fegray, anchor=south east] at (axis cs:30,420) {bad years first: ends at \$85,000};
\end{axis}
\end{tikzpicture}
\caption{Sequence-of-returns risk. Two retirees start with \$1 million and withdraw \$40,000 a year in real terms. Both experience the same thirty annual returns, three years of $-15\%$ and twenty-seven of $+7\%$, in opposite orders. With the bad years first, the portfolio falls to \$85,000; with the bad years last, it ends at \$1.77 million.}
\label{fig:03-sequence}
\end{figure}

\section{How it works}
\label{sec:03-how}

Investors can lose return to their own behavior through three channels.

\textbf{Timing.} Money flows into funds and markets after good performance and out after bad performance. If returns do not persist, or partly reverse, the average dollar is invested at relatively high prices. The gap between money-weighted and time-weighted returns measures the damage.

\textbf{Trading.} Every trade costs commissions, bid--ask spreads\index{bid--ask spread} and, in taxable accounts, taxes\index{taxes}. Trading would pay if investors could pick winners, but for most individuals the stocks bought do no better than the stocks sold. The costs are then a pure loss.

\textbf{Selling the wrong things.} The disposition effect leads investors to realize gains and defer losses. That raises taxes, because gains are taxed and losses are not harvested. It also costs return if past winners\index{winners and losers} tend to keep winning, as momentum research suggests (Chapter~\ref{ch:05}).

Not every gap is a mistake. Investors add money when they have it, often when incomes are high and markets have risen, and withdraw it when they need it. A money-weighted return below the time-weighted return can reflect the pattern of savings and spending, not bad judgment. Separating the two is the main difficulty in measuring the cost of behavior.

\section{The claim}
\label{sec:03-claim}

\begin{callout}
\textbf{Claim.} Individual investors lose a meaningful part of the return on their investments through their own trading behavior.
\end{callout}

The claim has three parts:

\begin{itemize}
  \item \textbf{A gap exists:} investors earn less than the investments they hold, after accounting for when they had money to invest.
  \item \textbf{Behavior causes it:} the gap comes from timing, overtrading and the disposition effect, not from costs the investor could not avoid.
  \item \textbf{It is meaningful:} large enough to matter over a lifetime, and present across markets and decades.
\end{itemize}

\section{The original evidence}
\label{sec:03-evidence}

\textbf{Overtrading (Barber and Odean 2000).} Barber and Odean studied 66,465 households with accounts at a large U.S. discount broker from 1991 to 1996. The average household turned over about 75\% of its portfolio a year and earned about 16.4\% a year, against 17.9\% for the market. The fifth of households that traded most earned about 11.4\%. The households did not pick bad stocks so much as pay too much to trade them.

\textbf{The disposition effect (Odean 1998).}\index{disposition effect|textbf} Using the trades of 10,000 accounts at a discount broker from 1987 to 1993, Odean found that investors realized gains much more readily than losses. Rebalancing and the higher trading costs\index{transaction costs} of low-priced stocks did not explain the pattern, and the subsequent performance of the stocks did not justify it. The effect reversed in December, when investors sold losers for tax reasons. Holding losers and selling winners lowered after-tax returns.

\textbf{Dollar-weighted market returns (Dichev 2007).} Dichev computed the dollar-weighted return of the aggregate investor in each stock market, using the flows into and out of equities implied by issuance and repurchases. Dollar-weighted returns were about 1.3 percentage points a year lower than buy-and-hold returns for NYSE and AMEX stocks from 1926 to 2002, about 5.3 points lower for Nasdaq from 1973 to 2002, and about 1.5 points lower across major international markets. Capital tended to flow in before low returns and out before high ones.

\textbf{Dumb money (Frazzini and Lamont 2008).} Frazzini and Lamont studied the flows of retail investors across U.S. mutual funds from 1980 to 2003. They used each fund's flows to measure how much individual investors wanted to own the stocks that fund held. Investors directed their money to funds holding stocks that went on to earn low returns. A short-lived ``smart money'' effect lasted about one quarter, but at longer horizons the ``dumb money'' effect dominated, so their reallocations lowered their wealth. The stocks they favoured also tended to be ones whose companies were issuing new shares, which links the fund-level evidence to Dichev's market-level result: companies sell shares when investors are keen to buy them.

\textbf{Fund investors (Morningstar 2024).} Morningstar's annual \emph{Mind the Gap} study compares the dollar-weighted returns of U.S. fund investors with the funds' total returns\index{return!total}. In the 10 years to December 2023, the average dollar invested earned 6.3\% a year, about 1.1 percentage points less than the funds, and investors trailed in every one of the 10 calendar years. The gap was smallest for allocation funds, which rebalance automatically, and largest for volatile sector funds.

\textbf{Other markets (Barber, Lee, Liu and Odean 2009).} Using complete trading records from Taiwan, Barber, Lee, Liu and Odean found that individual investors lost 3.8 percentage points a year to their trading, a total equal to 2.2\% of Taiwan's GDP. Virtually all the losses came from aggressive orders, while institutions profited.

\textbf{App-based trading (Barber, Huang, Odean and Schwarz 2022).} Robinhood users, trading through an app with zero commissions, concentrated their buying in stocks that attracted attention. The stocks they bought most heavily each day went on to earn abnormal returns\index{return!abnormal} of $-4.7\%$ over the following 20 days on average. Commission-free trading removed one cost but not the behavior.

The evidence for the first two parts of the claim is broad and consistent across countries and decades. The size of the loss, the third part, depends on how it is measured.

\section{Lab~\thechapter: Replicating Dichev (2007), Odean (1998) and Barber and Odean (2000)}
\label{sec:03-lab}

\begin{callout}
\textbf{Lab~\ref{ch:03} at a glance.} \emph{Claim on trial\index{claims-on-trial method}:} individual investors lose return to their own behavior. \emph{Papers replicated:} Dichev (2007), exactly, at the level of the U.S. market; the fund-investor gap measured by Dichev's method; Odean (1998) and Barber and Odean (2000), by simulation on real industry returns; Bengen (1994). \emph{Data, all free:} net equity issues and the market value of U.S. corporate equities from the corporate-equities tables of the Federal Reserve's \emph{Financial Accounts of the United States} (Z.1), annually from 1946; monthly net new cash flow to U.S. domestic equity mutual funds from the Investment Company Institute, through DataHub's public copy, which starts in 2007; the market return, the Fama--French\index{Fama--French factors} factors and 49 industry portfolios from the Kenneth R. French Data Library\index{Kenneth French data library}; Robert Shiller's U.S. stock market data from 1871. Individual brokerage records are private, which is why Parts C and D use simulated accounts. \emph{Code:} runs as a script (\texttt{python lab03.py -{}-z1-file ... -{}-flows-file ... -{}-shiller-file ...}) or as a notebook (\texttt{lab03\_walkthrough.ipynb}) with a prediction prompt before each step.
\end{callout}

The lab has five parts:

\begin{itemize}
  \item \textbf{Part A, Dichev's dollar-weighted market return.} From the Z.1 tables the lab builds, for each year since 1946, the market value of U.S. corporate equities and the net cash flow to shareholders of Equation~\eqref{eq:03-netflow}, with dividends taken from Shiller's data. It solves for the dollar-weighted return over the whole period and by 20-year window, and compares it with the buy-and-hold return on the French market portfolio. Dichev's estimate for NYSE and AMEX stocks, 1.3 percentage points a year from 1926 to 2002, is the benchmark the replication is judged against (Lab Figure~\ref{ch:03}.1).
  \item \textbf{Part B, the fund-investor gap.} Using monthly net flows to U.S. domestic equity mutual funds since 2007, and the market return as a proxy for the funds' return, the lab rebuilds the aggregate fund investor's assets. It compares the time-weighted return with the money-weighted return for the whole sample and for two halves, 2007--2015 and 2016 to the present. It also tests whether flows chase past 12-month returns and whether they predict the next 12 months (Lab Figure~\ref{ch:03}.2).
  \item \textbf{Part C, the disposition effect.} Four hundred simulated investors each hold 5 of the 49 French industry portfolios and sell one a month with 15\% probability. They differ only in how strongly they prefer to sell a winner when they hold both winners and losers. The lab reports Odean's PGR and PLR, and the next 12 months' return of the winners sold and the losers kept (Lab Figure~\ref{ch:03}.3).
  \item \textbf{Part D, overtrading.} Random portfolios of 10 industries replace holdings at random, at annual turnover\index{turnover} from 0\% to 250\%, paying 50 basis points\index{basis point} per trade. Because the trades are random, any difference in gross return is luck; the cost is the systematic effect, as in Barber and Odean (Lab Figure~\ref{ch:03}.4).
  \item \textbf{Part E, withdrawals.} Using Shiller's data from 1871, the lab tests a fixed inflation-adjusted withdrawal (Bengen's rule), a fixed percentage of current wealth and a simple guardrail rule. Each is applied to a 60/40 portfolio over every 30-year period, with the sequence-of-returns risk illustrated (Lab Figures~\ref{ch:03}.5 and~\ref{ch:03}.6).
\end{itemize}

The command \texttt{python lab03.py -{}-source sim} tests the code on simulated data.

\textbf{What the simulated run shows.} Parts B to E were run end to end on simulated data; Part A runs only on the Z.1 data, and its results are printed from the published run. In the simulation, fund investors chase past returns, industries have persistent return differences (so past winners tend to keep winning), and annual stock and bond returns have U.S.-like means and volatilities\index{volatility}. The results test the code and illustrate the mechanics.

\begin{itemize}
  \item Over 50 simulated years, the fund earned 11.7\% a year buy-and-hold while its investors earned 10.7\%, a gap of 1.0 percentage point. Flows followed past returns strongly ($t = 31$) but did not predict future returns. The gap came from money arriving after good years, not from foresight in reverse.
  \item Investors with no preference realized gains and losses at similar rates (PGR 0.19, PLR 0.21). Those with a strong preference for selling winners had a PGR of 0.42 and a PLR of 0.07, close to the pattern Odean found. The winners they sold went on to beat the losers they kept by 0.8 percentage points over the next year.
  \item Before costs, random portfolios earned about 4.2\% a year at every turnover level. After 50 basis points per trade, 75\% turnover cut the return to 3.4\% and 250\% turnover to 1.7\%, about one percentage point lost for each 100\% of turnover.
  \item A 4\% fixed real withdrawal from a 60/40 portfolio survived 96\% of 30-year periods; 5\% survived 74\% and 6\% only 46\%. Withdrawing a fixed percentage of current wealth never ran out, but in the worst period spending fell to about 38\% of its first-year level at a 4\% rate.
\end{itemize}

On real data, the lab measures Dichev's gap for the U.S. market since 1946, the gap for actual fund flows since 2007, the cost of the disposition effect when trading industries, and the historical record of withdrawal rules.

\emph{[Draft note: the Part A figures and all real-data results are generated by running the lab at the time of publication. The text above does not anticipate them.]}

\section{Stress tests}
\label{sec:03-stress}

\textbf{Measurement.} The size of the behavior gap depends on the method. Dollar-weighted returns are sensitive to large flows near the end of the sample and to the starting date. Comparing a dollar-cost-averaging investor with a lump-sum benchmark, as some industry studies appear to do, mixes the effect of behavior with the effect of when savings were available.

\textbf{Savings versus mistakes.} Investors often add money when incomes are high, which tends to be after strong markets, and withdraw it when they need it, sometimes in recessions. Part of any gap reflects the timing of savings and needs, not avoidable mistakes.

\textbf{Which investors.} Brokerage studies cover active traders at discount brokers, who may trade more than the typical saver. Most retirement savers contribute automatically and rarely trade, and their gap is likely smaller.

\textbf{New platforms.} Zero-commission apps removed explicit costs, but studies of app users find attention-driven buying that earns poor subsequent returns. Options trading by individuals, which has grown rapidly since 2020, adds costs through bid--ask spreads.

\textbf{Other countries and periods.} Taiwan's complete records show larger losses than U.S. brokerage samples. The pattern of performance chasing appears in fund flows in many countries.

\textbf{Withdrawal rules.} Historical success rates for withdrawal rules come from one country's unusually good century of returns. Results are less favourable using international data, lower expected returns\index{return!expected} or fees.

\section{Verdict}
\label{sec:03-verdict}

\begin{callout}
\textbf{Verdict: survives} (as of October 2026). Individual investors, on average, earn less than the investments they hold, and their behavior is a major reason. The size of the loss is smaller than the most quoted figures suggest.
\end{callout}

The evidence that investors hurt themselves:

\begin{itemize}
  \item Active traders lose return to trading costs (Barber and Odean 2000), and in complete-market data the losses are large (Barber, Lee, Liu and Odean 2009).
  \item Investors sell winners and hold losers (Odean 1998), a pattern found in many markets.
  \item Aggregate flows chase past returns, and dollar-weighted returns fall short of buy-and-hold returns in stock markets (Dichev 2007) and mutual funds (Morningstar 2024).
  \item Commission-free trading has not removed the problem (Barber, Huang, Odean and Schwarz 2022).
\end{itemize}

Careful measures of the fund-investor gap are about 1 to 1.5 percentage points a year, not the 5 to 6 points of some industry reports, and part of even the careful measure reflects the timing of savings rather than mistakes. A 1-point gap is still large: over 30 years it reduces ending wealth by roughly a quarter.

The verdict would \emph{weaken} if careful measures showed the gap disappearing for investors using automatic contributions and all-in-one funds, which would make it a problem of product design rather than of individual investors.

\section{What investors can use}
\label{sec:03-investors}

The remedies are rules that take decisions out of the moment:

\begin{itemize}
  \item \textbf{Automate contributions.} Investing a fixed amount on a schedule removes the temptation to time the market, and buys more shares when prices are low.
  \item \textbf{Rebalance on a calendar.}\index{rebalancing|textbf} Rebalancing once a year, or when weights drift past set bands, forces the investor to sell what has risen and buy what has fallen, the opposite of performance chasing. All-in-one allocation funds do this automatically, and Morningstar finds their investors have the smallest gaps.
  \item \textbf{Write an investment policy.} A one-page statement of the target allocation, the rebalancing rule and what will not trigger a change is a commitment device for bad markets.
  \item \textbf{Trade less.} For most individuals, the expected value of an extra trade is its cost.
  \item \textbf{Harvest losses, not gains.} In taxable accounts, realizing losses and deferring gains reverses the disposition effect and lowers taxes.
  \item \textbf{Diversify broadly.}\index{diversification|textbf} Broad index funds remove the temptation and the risk of concentrated bets on familiar or attention-grabbing stocks.
\end{itemize}

\textbf{The lifecycle extension: spending in retirement.} Behavior matters most when an investor starts drawing on savings.

\begin{itemize}
  \item Bengen's 4\% rule is a reasonable starting point for a 30-year retirement with a balanced portfolio, based on U.S. history. It is not a guarantee: it depends on future returns resembling the past, and it ignores fees.
  \item Fixed real withdrawals carry sequence-of-returns risk: a bad decade early in retirement can exhaust a portfolio that would have lasted had the same returns come in a different order.
  \item Flexible rules trade certainty of income for certainty of not running out. A fixed percentage of current wealth never runs out but can force deep spending cuts. \textbf{Guardrail rules} adjust spending only when it drifts far from the plan.
  \item Annuities, which insure against outliving one's savings, are the main alternative for the part of spending that must not fall.
\end{itemize}

\textbf{Questions to ask an adviser or a fund provider:}

\begin{itemize}
  \item What return did investors in this fund actually earn (the money-weighted return), and how does it compare with the fund's reported return?
  \item Is rebalancing automatic?
  \item What does the plan say to do after a 30\% fall in the market?
  \item How is the retirement withdrawal rate\index{withdrawal rate} set, and what happens to spending if the first five years are bad?
\end{itemize}

\begin{callout}
\textbf{What Failed: the 6-point behavior gap.} For years, financial firms and advisers quoted DALBAR's annual \emph{Quantitative Analysis of Investor Behavior} as proof of how badly investors behave. In its 2016 edition, covering the 20 years to 2015, DALBAR reported that the S\&P 500 returned 10.35\% a year while the average fund investor earned 3.66\%, a gap of more than six percentage points (DALBAR 2016). Wade Pfau and others argued that the methodology, which DALBAR does not fully disclose, appeared to compare a lump-sum investment in the index with investors who added money over time, mixing the effect of behavior with the timing of savings. Using fund-level flows from 1990 to 2016, David Blanchett of Morningstar Investment Management found a gap of about 1.6 percentage points a year for equity funds and about 1.5 points on average across asset classes (Blanchett 2017). The claim that investors hurt themselves survived; the dramatic number did not.
\end{callout}

\section{Open questions}
\label{sec:03-open}

\begin{itemize}
  \item \textbf{Savings or mistakes?} How much of the measured behavior gap reflects the unavoidable timing of savings and withdrawals rather than avoidable errors?
  \item \textbf{Default design.} Do target-date funds, automatic enrolment and automatic rebalancing close the gap, or do investors find new ways to trade?
  \item \textbf{New products.} What do zero-day options, leveraged ETFs and cryptocurrency trading cost individual investors, and are the losses concentrated in a small group?
  \item \textbf{Advice.} Do human advisers, robo-advisers or language-model assistants reduce behavioral losses enough to cover their fees?
\end{itemize}

\section*{Exercises}
\addcontentsline{toc}{section}{Exercises}

\subsection*{Conceptual (all readers)}

\begin{enumerate}
  \item A fund returns $+30\%$ in year 1 and $-20\%$ in year 2. Investor A invests \$1,000 at the start. Investor B invests \$1,000 at the start and \$3,000 at the start of year 2. Compute the fund's time-weighted return and each investor's money-weighted return.
  \item On a day an investor sells, she holds 4 stocks at a gain and 6 at a loss and sells 2 winners and 1 loser. Compute that day's contributions to realized and paper gains and losses, and her PGR and PLR for the day.
  \item A retiree starts with \$1 million and withdraws \$40,000 a year, adjusted for inflation. Explain why two retirees with the same average real return over 30 years can end with very different wealth.
\end{enumerate}

\subsection*{Lab extensions (technical readers)}

\begin{enumerate}
  \setcounter{enumi}{3}
  \item Part A: compare your dollar-weighted market return since 1946 with Dichev's 1.3 percentage points for 1926--2002. How much does the answer change if you start in 1980, the start of the buyback era?
  \item Part B: run on the ICI flow data. What is the behavior gap for the whole sample and for each half? Do flows chase past returns, and do they predict the next 12 months?
  \item Part C: compare the cost of the disposition effect with that in the simulation. How does industry momentum in the real data change the result?
  \item Part D: how much annual turnover can an investor afford before costs of 20 basis points per trade cost 0.5\% a year?
  \item Part E: find the five worst start years for a 4\% fixed real withdrawal in Shiller's data. What happened in the first decade of each?
  \item A firm advertises that ``the average investor earns 6\% a year less than the market''. What questions would you ask about how the figure was computed?
  \item Design a default retirement plan for employees who will never read their statements. Which choices from this chapter would you build in, and why?
\end{enumerate}
